%!TEX root = ../../main.tex

\chapter{Grundlagen}
\label{ch:grundlagen}

\section{ETL-Prozesse}
\label{sec:etl}

\subsection{Definition und Phasen}
\label{subsec:etl-phasen}
% Extract, Transform, Load - Definitionen und typischer Ablauf.

\subsection{Typische Herausforderungen}
\label{subsec:etl-herausforderungen}
% Datenqualität, Heterogenität, Skalierbarkeit, Wartbarkeit.

\subsection{Klassische Werkzeuge und Verfahren}
\label{subsec:etl-werkzeuge}
% z.B. Pandas, Apache Spark, Talend, Pentaho, regelbasierte Ansätze.

\section{Datenqualität und Datenbereinigung}
\label{sec:datenqualitaet}

\subsection{Dimensionen der Datenqualität}
\label{subsec:dq-dimensionen}
% Genauigkeit, Vollständigkeit, Konsistenz, Aktualität, Eindeutigkeit.

\subsection{Häufige Datenprobleme}
\label{subsec:dq-probleme}
% Fehlende Werte, Duplikate, inkonsistente Formate, Ausreißer.

\section{Large Language Models}
\label{sec:llm}

\subsection{Funktionsweise und Architektur}
\label{subsec:llm-architektur}
% Transformer, Tokenisierung, Training, Inferenz.

\subsection{Überblick aktueller Modelle}
\label{subsec:llm-modelle}
% GPT-Familie, Claude, Gemini, Llama, etc.

\subsection{Stärken und Limitationen bei strukturierten Daten}
\label{subsec:llm-strukturierte-daten}
% Halluzinationen, Kontextfenster, Kosten, Reproduzierbarkeit.

\section{Prompt Engineering}
\label{sec:prompting}

\subsection{Grundlegende Prompting-Strategien}
\label{subsec:prompting-strategien}
% Zero-Shot, Few-Shot, Chain-of-Thought, Role-Prompting.

\subsection{Einflussfaktoren auf die Ergebnisqualität}
\label{subsec:prompting-einfluss}
% Kontextlänge, Beispielwahl, Temperatur, Modellgröße.
